\documentclass[runningheads]{llncs}
\usepackage[T1]{fontenc}
\usepackage{graphicx}
\usepackage{booktabs}
\usepackage{amsmath}
\usepackage{url}
\usepackage[hidelinks]{hyperref}

\begin{document}

\title{Public Ensemble Weights Outperform Six Retraining Attempts
in Interactive Whole-Body Lesion Segmentation}
\subtitle{A Negative-Results Report from a Single-GPU Participant in autoPET V}

\titlerunning{Public Weights vs.\ Retraining in autoPET V}

\author{FIRST AUTHOR\inst{1}\orcidID{0000-0000-0000-0000}}
\authorrunning{F. Author}
\institute{Hirosaki University, Hirosaki, Aomori, Japan\\
\email{AUTHOR@EXAMPLE.AC.JP}}

\maketitle

\begin{abstract}
We report our participation in the autoPET V challenge on interactive lesion
segmentation in whole-body FDG- and PSMA-PET/CT. Working on a single consumer
GPU (RTX 3060, 12\,GB), we retrained the challenge baseline six times under
controlled changes: a repaired click-channel pipeline, an on-the-fly click
trainer, a click-dropout trainer, a residual-encoder architecture, a two-model
ensemble, and three inference-side variants. None of them reached the
organizers' own baseline on the preliminary leaderboard. On the fold-0
validation split, evaluated with one evaluator on identical data, the official
baseline scores $0.7556$ against our best retrained model at $0.7733$; once the
number of cases entering each mean is equalized, the official baseline is
$+0.041$ ahead, and our apparent advantage is an artifact of $16$ lesion-free
cases being excluded from our mean but not from theirs. We therefore submit a
$10$-fold ensemble of the publicly released LesionLocator weights
(Apache-2.0), which the challenge rules explicitly permit. Beyond the
submission, we document three measurement failures that cost us most of the
challenge period: the preliminary test set contains five cases of which four
are scored, so leaderboard differences below roughly $0.05$ are noise; training
pseudo-Dice, single-case Dice, and a local multi-case interaction simulator each
failed to predict leaderboard order; and a genuine click-blindness bug turned
out not to explain the gap it was blamed for.

\keywords{autoPET challenge \and Interactive segmentation \and PET/CT \and
Negative results \and Reproducibility}
\end{abstract}

\section{Introduction}

The autoPET V challenge (MICCAI 2026) targets interactive lesion segmentation
in whole-body PET/CT. Unlike the fully automatic tasks of earlier autoPET
editions, an algorithm is called repeatedly: at the first iteration it receives
the CT and PET volumes with an empty scribble set, and at each subsequent
iteration it receives one additional scribble derived from its own previous
error, namely a foreground scribble on an under-segmented region or a
background scribble on an over-segmented one, whichever component is smaller.
The evaluation loop runs five iterations per case and ranks submissions by the
area under the Dice curve and the area under the detection-F1 (DMM) curve,
weighted equally.

Two properties of this protocol drive everything that follows. First, the
zero-click iteration is included in the area, so a model that is strong only
when clicked is heavily penalized. Second, each iteration adds exactly one
scribble, so by the last iteration the model has seen at most four. The regime
is sparse interaction, not dense correction.

We participated with a single consumer GPU. This report is deliberately
weighted toward what did not work: our six retrained models, the diagnostics
that identified and repaired a real defect in our training data, and the
measurement mistakes that let us pursue a wrong hypothesis for several weeks.
Our final submission uses publicly released weights rather than our own
training, and we state plainly why.

\section{Methods}

\subsection{Data}

\paragraph{Training data.}
We used only the data provided by the challenge: 1611 whole-body PET/CT studies
(FDG and PSMA tracers) with lesion annotations and simulated scribbles. No
external or private data were used for the models reported here. All training
used nnU-Net fold 0 of the organizers' \texttt{splits\_final.json}
(1288 training / 323 validation cases), which we verified reproduces the
organizers' own split: applying nnU-Net's rule (sorting the case keys, then
\texttt{KFold(5, shuffle=True, random\_state=12345)}) to the 1611 case
identifiers yields exactly the 1288/323 partition reported in the official
training log.

\paragraph{Validation data.}
Fold 0 of the same split (323 cases). We did not perform a full 5-fold
cross-validation: a single 1000-epoch run takes approximately three days on our
hardware, and the challenge period did not allow five of them. All validation
numbers below are therefore fold-0 numbers, and we report the number of cases
entering every mean, for reasons that Section~\ref{sec:results} makes clear.

\subsection{Data pre-processing}

The network input has four channels: CT, PET, a foreground-click map and a
background-click map. Because our early models were effectively blind to the
click channels, we reverse-engineered the organizers' click encoding rather
than assuming it. Regenerating the click maps from the released scribble JSON
files and comparing them against the organizers' pre-computed heatmaps gave a
byte-exact match on all 1611 cases: 1038 cases with clicks (2076 files, zero
mismatches, zero out-of-bounds coordinates, no values other than $1.0$) and 573
cases whose scribble lists are empty, for which the official heatmaps are
uniformly zero.

The encoding is therefore a single voxel set to $1.0$ with no smoothing, with
the JSON coordinate triple indexed directly into the NIfTI array in
\texttt{[x,y,z]} order. Normalization follows the official configuration:
CT normalization for the CT channel and Z-score normalization for the remaining
three, including the sparse binary click channels. We note that Z-score
normalization of a channel whose standard deviation is of order $10^{-4}$
amplifies the click voxels by roughly three orders of magnitude relative to the
image channels, whereas substituting no-normalization, as we did for our
on-the-fly models, leaves the clicks at $1.0$ where they are dominated by the
image channels. This is a plausible contributor to the weak click response we
measured in those models, although we did not isolate it experimentally.

\subsection{Algorithm/model}

\paragraph{Final submission.}
Our final submission is a 10-fold ensemble of the publicly released
LesionLocator interactive weights~\cite{lesionlocator}, distributed under the
Apache-2.0 license by the team that won the corresponding autoPET IV task. The
architecture is a residual-encoder U-Net (ResEnc-L) with a $192^3$ patch and
Euclidean-distance-transform click channels. The challenge rules state that
publicly available data and networks may all be used, and we make no claim of
novelty for the network itself.

Two changes were required to run the released code inside the challenge
container: extending the fold list to all ten folds, and disabling test-time
mirroring. We additionally strip each checkpoint to the four keys the nnU-Net
predictor actually reads (\texttt{network\_weights}, \texttt{trainer\_name},
\texttt{inference\_allowed\_mirroring\_axes}, \texttt{init\_args}), which
reduces each fold from $820$\,MB to $391$\,MB and the ten-fold model directory
from $8.2$\,GB to $3.9$\,GB. The stripped checkpoints are reloaded and verified
before the image is built.

\paragraph{Models we trained ourselves.}
For the negative results we report six configurations, all derived from the
challenge baseline (plain convolutional U-Net, six stages, patch
$112\times160\times128$):
(i) a retrained baseline after repairing the click channels;
(ii) an on-the-fly click trainer that synthesizes EDT click channels from the
ground truth at every iteration;
(iii) a click-dropout trainer ($p_{\mathrm{drop}}=0.5$) intended to break the
shortcut described below;
(iv) the on-the-fly trainer on a residual-encoder (ResEnc-M) backbone;
(v) a softmax-average ensemble of (ii) and (iv);
and (vi) three inference-side variants of the official weights
(\texttt{checkpoint\_best}, halved sliding-window step, full test-time
mirroring).

\subsection{Data post-processing}

None. The predicted label map is written at the original geometry with no
connected-component filtering, no size threshold and no morphological
operation.

\subsection{Training and test parameters}

Our own training used the nnU-Net defaults that the organizers also used, which
we read from the baseline's \texttt{debug.json} rather than assuming: SGD with
learning rate $0.01$, Nesterov momentum $0.99$, weight decay $3\times10^{-5}$,
polynomial learning-rate decay, 1000 epochs of 250 iterations, batch size 2,
patch $112\times160\times128$, foreground oversampling $0.33$, and the standard
nnU-Net spatial and intensity augmentations. A single run took approximately
three days on an RTX 3060.

At test time the final submission uses a sliding-window step size of $0.5$ with
mirroring disabled. One iteration of the 10-fold ensemble takes $171.9$\,s on
our RTX 3060 ($161.0$\,s of prediction, about $12.9$\,s per fold, plus roughly
$11$\,s of container startup), that is $859.5$\,s for the five iterations of one
case against the 15-minute budget. On the challenge hardware, which uses a
faster GPU, we expect this to be shorter, but we did not measure it.

\subsection{Github repository}

Link to Github repository: \url{https://github.com/leaieaie/autoPETV}

\section{Results}
\label{sec:results}

\paragraph{Fold-0 validation.}
Table~\ref{tab:val} compares the official baseline weights with our best
retrained model on the same 323-case validation split, using one evaluator, one
nnU-Net version and one prediction path. The comparison is not the one we first
believed. nnU-Net's mean validation Dice excludes cases in which the union of
prediction and ground truth is empty. Our click-dependent model predicts
nothing on lesion-free cases and is excluded from all 116 of them; the official
baseline produces a false positive on 16 of those cases and is therefore scored
$0.0$ on them. Equalizing the population reverses the ranking: restricted to
the 207 cases with lesions, the official baseline is ahead by $0.041$.

\begin{table}[t]
\caption{Fold-0 validation Dice, identical data and evaluator. Here $n$ is the
number of cases entering the mean. The reported means are not comparable; the
right-hand column is.}
\label{tab:val}
\centering
\begin{tabular}{lccc}
\toprule
Model & $n$ & Reported mean & Lesion-bearing only ($n=207$) \\
\midrule
Official baseline & 223 & 0.7556 & \textbf{0.8140} \\
Ours (repaired clicks) & 207 & 0.7733 & 0.7733 \\
\bottomrule
\end{tabular}
\end{table}

We also note that the mean validation Dice recorded in the organizers' own
training log ($0.7225$) differs from our re-measurement of the same weights
($0.7556$) by $+0.033$. Re-running our own model from raw data rather than from
preprocessed data changed its score by $2.5\times10^{-7}$, so the discrepancy is
not the prediction path; we attribute it to library or environment versions.
Numbers taken from the official log should not be compared against numbers
produced elsewhere.

\paragraph{Preliminary leaderboard.}
Table~\ref{tab:lb} lists every submission we made. The official baseline, shown
for reference, was not beaten by any of them.

\begin{table}[t]
\caption{Preliminary test phase, AUC-Dice and AUC-F1. All values are means over
four scored cases (see Section~\ref{sec:discussion}).}
\label{tab:lb}
\centering
\begin{tabular}{lcc}
\toprule
Submission & AUC-Dice & AUC-F1 \\
\midrule
Official baseline (reference) & \textbf{0.7800} & \textbf{0.7554} \\
\midrule
On-the-fly click trainer & 0.7525 & 0.6622 \\
Click dropout ($p=0.5$) & 0.7161 & 0.6724 \\
ResEnc-M on-the-fly & 0.6805 & 0.7059 \\
Two-model ensemble & 0.7189 & 0.7263 \\
Official weights, \texttt{checkpoint\_best} & 0.7709 & 0.7545 \\
Official weights, step size $0.25$ & 0.7781 & 0.7431 \\
Official weights, full mirroring & 0.7826 & 0.7461 \\
LesionLocator, 6 folds & 0.7334 & 0.6030 \\
\bottomrule
\end{tabular}
\end{table}

\paragraph{Interaction behaviour.}
A six-step diagnostic on individual cases separates two distinct failure modes.
Our first repaired model responds strongly to clicks, with the predicted volume
rising from 847 to 20431 voxels over the iterations and a final-iteration DMM of
$0.842$, but it collapses at the zero-click iteration (Dice $0.071$). The cause
is a shortcut that is genuinely present in the data: of the 1611 training
cases, 573 have empty scribble lists, and in a 300-case audit the
correspondence between ``no tumor clicks'' and ``empty label'' was exact, with
189 cases with tumor clicks all non-empty and 111 cases without all empty. A
model can satisfy the training objective by predicting nothing when no click is
given. Click dropout repairs the zero-click iteration (Dice $0.071$ to
$0.861$) but weakens the click response, and scored lower on the leaderboard.

\section{Discussion}
\label{sec:discussion}

\paragraph{The preliminary phase cannot rank anything.}
The preliminary test set contains five cases, one of which is lesion-free and
scored as null, so every leaderboard figure in Table~\ref{tab:lb} is a mean
over four cases. Per-case metrics are downloadable from the evaluation page,
and we should have looked at them on our first submission rather than our
fifteenth. Testing the per-case differences between our best public-weight run
and the official baseline gives $t=-0.39$: the $0.047$ gap that we spent days
interpreting is not distinguishable from zero. The same holds for the
$\pm0.005$ differences among our three inference-side variants, which consumed
three submission slots and three days. We report this because the ranking on a
four-case leaderboard is visible to every participant and invites exactly the
over-interpretation we fell into.

\paragraph{Three predictors that do not work.}
We tried, in order, to select models by training pseudo-Dice, by single-case
Dice, and by a local multi-case interaction simulator. Each was contradicted by
the leaderboard. Pseudo-Dice preferred an epoch-992 checkpoint that scored
$0.009$ lower than the epoch-1000 checkpoint. Single-case Dice predicted the
wrong direction for both inference variants we checked. The local simulator
ranked our click-dropout model above our on-the-fly model; the leaderboard
ranked them the other way. We do not have a working local proxy for this
challenge's metric, and we think this is the main reason our search was
unproductive: without one, every hypothesis costs a submission slot.

\paragraph{A real bug that was not the cause.}
Our preprocessed click channels were entirely zero for every training case, a
defect introduced by a data-preparation script that reported no errors. This is
a genuine bug, it is fully documented, and repairing it changed the mean
validation Dice from $0.4371$ to $0.7733$. It is nevertheless not the
explanation for our leaderboard deficit, for two reasons: our on-the-fly
trainer synthesizes its click channels during training and never reads the
broken ones, yet still lost to the baseline; and every model trained after the
repair also lost. Having found one sufficient-looking cause, we stopped looking
and wrote it up as settled. The actual reason our retrained models trail the
official baseline, trained on the same data, the same split, the same
architecture and the same schedule, remains unidentified. Candidates we have
not tested include library-version differences and the normalization of the
click channels discussed above.

\paragraph{Why the area is dominated by the first iteration.}
Because the loop starts with no scribbles and adds one per iteration, AUC-Dice
is largely a measure of fully automatic segmentation quality. A model tuned for
responsiveness to interaction, as ours were, optimizes the part of the curve
that carries the least weight. In hindsight the click-dropout direction, namely
forcing the model to segment well with no clicks at all, was the right one, and
$p_{\mathrm{drop}}=0.5$ was simply too aggressive.

\paragraph{On submitting public weights.}
Our final submission is not our own model. The challenge rules permit publicly
available networks without qualification, and we prefer to state the situation
accurately: on this task, with this data, a single consumer GPU and roughly
three months, we did not produce a model that beats either the organizers'
baseline or the released weights of the previous edition's winner. The
contribution we can defend is the measurement record above, not the network.

\section{Conclusion}

We submit a 10-fold ensemble of publicly released Apache-2.0 weights for
autoPET V, and report six retraining attempts that failed to reach the official
baseline. The most transferable lessons are methodological: verify the encoding
of auxiliary input channels against the organizers' own artifacts before
training on them; report the number of cases entering every mean, because
nnU-Net's exclusion of empty cases can reverse a comparison; and check how many
cases a preliminary leaderboard actually scores before treating its ordering as
information.

\begin{table}[t]
\caption{Algorithm details.}
\label{tab:details}
\centering
\begin{tabular}{p{2.2cm}p{3.0cm}p{3.2cm}p{2.4cm}}
\toprule
Team name & Algorithm name (as submitted on grand-challenge) &
Data pre-processing & Data post-processing \\
\midrule
TEAM NAME &
nnUNet 3d fullres interactive (CT/PET + scribbles) &
nnU-Net default resampling; CT normalization for CT, Z-score for PET and both
click channels; EDT click encoding &
None \\
\bottomrule
\end{tabular}
\end{table}

\begin{credits}
\subsubsection{\ackname} We thank the autoPET organizers for releasing the
baseline training artifacts, without which the analysis in
Section~\ref{sec:results} would not have been possible.

\subsubsection{\discintname}
The author has no competing interests to declare that are relevant to the
content of this article.
\end{credits}

\begin{thebibliography}{8}

\bibitem{lesionlocator}
Rokuss, M., et al.: LesionLocator: Zero-shot universal tumor segmentation and
tracking in 3D whole-body imaging. arXiv:2508.21680 (2025)

\bibitem{nnunet}
Isensee, F., Jaeger, P.F., Kohl, S.A.A., Petersen, J., Maier-Hein, K.H.:
nnU-Net: a self-configuring method for deep learning-based biomedical image
segmentation. Nature Methods \textbf{18}(2), 203--211 (2021)

\bibitem{autopetdata}
Gatidis, S., et al.: A whole-body FDG-PET/CT dataset with manually annotated
tumor lesions. Scientific Data \textbf{9}, 601 (2022)

\bibitem{resenc}
Isensee, F., et al.: nnU-Net revisited: a call for rigorous validation in 3D
medical image segmentation. In: MICCAI 2024, LNCS. Springer (2024)

\end{thebibliography}

\end{document}
